% =========================================================================
% ACTA DE REUNIÓN — CS3081-007-2026
% Plantilla de estilo basada en el modelo oficial de Google Docs del curso
% Compilar con:  tectonic -X compile CS3081-007-2026_Acta_Reunion7.tex
% =========================================================================
\documentclass[11pt,a4paper]{article}

\usepackage[a4paper, top=2.6cm, bottom=2.4cm, left=2.5cm, right=2.5cm]{geometry}
\usepackage{fontspec}
\setmainfont{Carlito}                    % métrica idéntica a Calibri
\usepackage[table]{xcolor}
\usepackage{array}
\usepackage{longtable}
\usepackage{fancyhdr}
\usepackage{lastpage}
\usepackage{needspace}

% ------------------------- Paleta de la plantilla -------------------------
\definecolor{rojo}{HTML}{E4002B}   % marca CS3081
\definecolor{grisB}{HTML}{B7B7B7}  % bordes de tabla
\definecolor{grisH}{HTML}{202124}  % encabezados de tabla (negro)
\definecolor{grisL}{HTML}{F3F4F6}  % celdas de etiqueta
\definecolor{rosa}{HTML}{FCE8EC}   % caja PROPÓSITO / REGISTRO
\definecolor{grisV}{HTML}{555555}  % valores (itálica)
\definecolor{grisN}{HTML}{666666}  % notas y pie
\definecolor{tinto}{HTML}{333333}

\arrayrulecolor{grisB}
\renewcommand{\arraystretch}{1.55}
\setlength{\tabcolsep}{4pt}

% ------------------------- Encabezado y pie corridos ----------------------
\pagestyle{fancy}
\fancyhf{}
\fancyhead[L]{\footnotesize\textbf{\color{rojo}CS3081\hspace{0.7em}|\hspace{0.7em}INGENIERÍA DE SOFTWARE}\hspace{1.4em}{\color{grisN}·\hspace{0.7em}Acta de reunión con usuarios}}
\fancyfoot[R]{\footnotesize\color{grisN}Uso académico · Proyecto real\hspace{2em}Página \thepage\ de \pageref{LastPage}}
\renewcommand{\headrulewidth}{0pt}
\renewcommand{\footrulewidth}{0pt}
\setlength{\headheight}{14pt}

% ------------------------- Comandos de estilo -----------------------------
\newcommand{\asec}[1]{\Needspace*{13\baselineskip}\par\vspace{18pt}{\noindent\fontsize{14}{17}\selectfont\textbf{#1}}\par\vspace{9pt}}
\newcommand{\cajanota}[2]{\par\noindent\fcolorbox{grisB}{rosa}{\parbox{\dimexpr\linewidth-2\fboxsep-2\fboxrule}{\footnotesize \textbf{\textcolor{rojo}{#1: }}{\color{tinto}#2}}}\par\vspace{7pt}}
\newcommand{\hc}[1]{\cellcolor{grisH}\textcolor{white}{\textbf{#1}}}
\newcommand{\lb}[1]{\cellcolor{grisL}\textbf{#1}}
\newcommand{\vl}[1]{\textcolor{grisV}{\textit{#1}}}
\newcolumntype{L}[1]{>{\raggedright\arraybackslash}p{#1}}
\newcolumntype{C}[1]{>{\centering\arraybackslash}p{#1}}

\setlength{\parindent}{0pt}
\setlength{\parskip}{4pt}

\begin{document}

% ============================ TÍTULO ======================================
{\centering {\fontsize{22}{26}\selectfont\textbf{ACTA DE REUNION}}\par}
\vspace{10pt}

% ============================ 1. DATOS ====================================
\asec{1. DATOS GENERALES}
{\footnotesize
\begin{tabular}{|L{2.35cm}|L{5.45cm}|L{2.1cm}|L{4.3cm}|}
\hline
\lb{Código de acta} & \vl{CS3081-007-2026} & \lb{Fecha} & \vl{21/09/2026} \\ \hline
\lb{Hora de inicio} & \vl{10:00 a.m.} & \lb{Hora de término} & \vl{10:30 a.m.} \\ \hline
\lb{Proyecto} & \multicolumn{3}{L{12.0cm}|}{\vl{Proyecto Igualab --- Plataforma de analítica de sostenibilidad (RAG)}} \\ \hline
\lb{Módulo/Fase} & \vl{Avance desplegado en los entornos de la universidad y demostración del \textbf{módulo de autenticación}; configuración del envío de correos con \textbf{Gmail} para el flujo de recuperación de contraseña; y \textbf{presentación del Análisis y Diseño (A\&D) por correo}} & \lb{Modalidad} & \vl{[ ] Presencial \ [X] Virtual} \\ \hline
\end{tabular}}

% ============================ 2. OBJETIVO =================================
\asec{2. OBJETIVO Y RESULTADO ESPERADO}
{\footnotesize
\begin{tabular}{|L{2.9cm}|L{11.3cm}|}
\hline
\lb{Objetivo de la reunión} & \vl{Revisar el \textbf{avance ya desplegado en los entornos de la universidad} y la demostración del \textbf{módulo de autenticación}; coordinar la \textbf{configuración de envío de correos con Gmail} (SMTP/API) para el flujo de \textbf{recuperación de contraseña}; y acordar la \textbf{presentación del Análisis y Diseño (A\&D) por correo}.} \\ \hline
\lb{Resultado esperado} & \vl{\textbf{Avance desplegado y módulo de autenticación mostrados} al PO; \textbf{acuerdo sobre la configuración de correo con Gmail} y \textbf{cuenta a compartir por el PO} para dicha configuración; y \textbf{Análisis y Diseño (A\&D) presentado/remitido por correo}.} \\ \hline
\end{tabular}}

% ============================ 3. PARTICIPANTES ============================
\newpage
\asec{3. PARTICIPANTES}
{\footnotesize
\begin{tabular}{|C{0.8cm}|L{4.1cm}|L{3.1cm}|L{4.4cm}|C{1.8cm}|}
\hline
\hc{N°} & \hc{Nombre y apellido} & \hc{Organización / equipo} & \hc{Rol en la reunión} & \hc{Asistencia} \\ \hline
1 & \vl{Oscar Rafael Baldeón Montoro} & \vl{Cliente} & \vl{Cliente / Product Owner} & \vl{[X] Sí \ [ ] No} \\ \hline
2 & \vl{Fabricio Godofredo Ladera La Torre} & \vl{Equipo UTEC} & \vl{Jefe de Proyecto (PM)} & \vl{[X] Sí \ [ ] No} \\ \hline
3 & \vl{Juan Renato Flores Pascual} & \vl{Equipo UTEC} & \vl{Desarrollador Frontend} & \vl{[ ] Sí \ [X] No} \\ \hline
4 & \vl{Carlos David Ordinola Ortega} & \vl{Equipo UTEC} & \vl{Analista Funcional} & \vl{[ ] Sí \ [X] No} \\ \hline
5 & \vl{Luis Javier Millones Carrasco} & \vl{Equipo UTEC} & \vl{Analista Funcional} & \vl{[X] Sí \ [ ] No} \\ \hline
6 & \vl{Juan Marcelo Ferreyra Gonzales} & \vl{Equipo UTEC} & \vl{Desarrollador Backend} & \vl{[X] Sí \ [ ] No} \\ \hline
7 & \vl{Mauricio Gabriel Gonzalez Kremer} & \vl{Equipo UTEC} & \vl{Tester} & \vl{[ ] Sí \ [X] No} \\ \hline
8 & \vl{Alonso Aarón Benites Camacho} & \vl{Equipo UTEC} & \vl{Tester} & \vl{[X] Sí \ [ ] No} \\ \hline
\end{tabular}}

% ============================ 4. AGENDA ===================================
\asec{4. AGENDA}
\vspace{4pt}
{\footnotesize
\begin{tabular}{|C{0.8cm}|L{8.6cm}|L{2.6cm}|C{2.2cm}|}
\hline
\hc{N°} & \hc{Tema} & \hc{Responsable} & \hc{Tratado} \\ \hline
1 & \vl{\textbf{Avance desplegado en los entornos de la universidad} y demostración del \textbf{módulo de autenticación}} & \vl{Juan Marcelo (Backend) / Equipo} & \vl{[X] Sí \ [ ] No} \\ \hline
2 & \vl{\textbf{Configuración de envío de correos con Gmail} (SMTP/API) para recuperación de contraseña; cuenta a compartir por el PO} & \vl{Oscar / Fabricio (PM)} & \vl{[X] Sí \ [ ] No} \\ \hline
3 & \vl{\textbf{Presentación del Análisis y Diseño (A\&D)} por correo} & \vl{Luis / Carlos (Analistas)} & \vl{[X] Sí \ [ ] No} \\ \hline
\end{tabular}}

% ============================ 5. DESARROLLO ===============================
\newpage
\asec{5. DESARROLLO DE LA REUNIÓN}
\vspace{4pt}
{\footnotesize
\begin{longtable}{|C{0.8cm}|L{3.0cm}|L{7.6cm}|L{2.8cm}|}
\hline
\rowcolor{grisH}\hc{N°} & \hc{Tema / necesidad} & \hc{Síntesis de lo conversado} & \hc{Conclusión / estado} \\ \hline
\endfirsthead
\hline
\rowcolor{grisH}\hc{N°} & \hc{Tema / necesidad} & \hc{Síntesis de lo conversado} & \hc{Conclusión / estado} \\ \hline
\endhead
1 & \vl{Avance desplegado y módulo de autenticación} & \vl{El equipo informa que el avance del sistema \textbf{ya se encuentra desplegado en los entornos de la universidad}. Se muestra al PO el \textbf{módulo de autenticación} (inicio de sesión y control de acceso por roles), evidenciando el funcionamiento del flujo sobre el entorno desplegado.} & \vl{\textbf{Desplegado y mostrado}} \\ \hline
2 & \vl{Configuración de correo (Gmail)} & \vl{Se conversa sobre el \textbf{envío de correos con Gmail}. Se requiere \textbf{conseguir la API/configuración} de una cuenta Gmail para el envío de correos del sistema, por ejemplo en el caso de \textbf{``olvidé contraseña''} (recuperación de acceso). El PO se compromete a \textbf{compartir una cuenta de correo} para realizar la configuración con ese correo.} & \vl{\textbf{Pendiente de cuenta del PO --- por configurar}} \\ \hline
3 & \vl{Presentación del Análisis y Diseño (A\&D)} & \vl{Se acuerda que el \textbf{Análisis y Diseño (A\&D)} se \textbf{presentará/remitirá por correo} al cliente, junto con las coordinaciones pendientes de la reunión.} & \vl{\textbf{Acordado --- envío por correo}} \\ \hline
\end{longtable}}

% ============================ 6. REQUISITOS ===============================
\asec{6. NECESIDADES, REQUISITOS Y CAMBIOS IDENTIFICADOS}
\vspace{5pt}
{\footnotesize
\begin{longtable}{|C{1.3cm}|L{5.1cm}|C{1.5cm}|C{1.3cm}|L{3.4cm}|C{1.6cm}|}
\hline
\rowcolor{grisH}\hc{ID} & \hc{Necesidad / requisito / cambio} & \hc{Tipo} & \hc{Prioridad} & \hc{Criterio o evidencia} & \hc{Estado} \\ \hline
\endfirsthead
\hline
\rowcolor{grisH}\hc{ID} & \hc{Necesidad / requisito / cambio} & \hc{Tipo} & \hc{Prioridad} & \hc{Criterio o evidencia} & \hc{Estado} \\ \hline
\endhead
REQ-24 & \vl{\textbf{Despliegue del avance en los entornos de la universidad}: el avance del sistema se encuentra desplegado en el entorno de desarrollo de la universidad.} & \vl{Demo} & \vl{Alta} & \vl{Entorno de la universidad operativo con el avance desplegado} & \vl{\textbf{Registrado}} \\ \hline
REQ-25 & \vl{\textbf{Demostración del módulo de autenticación}: se mostró al PO el módulo de autenticación (inicio de sesión y control de acceso por roles) sobre el entorno desplegado.} & \vl{Demo} & \vl{Alta} & \vl{Módulo de autenticación mostrado en la reunión} & \vl{\textbf{Registrado}} \\ \hline
REQ-26 & \vl{\textbf{Configuración de envío de correos con Gmail} (SMTP/API) para el flujo de \textbf{recuperación de contraseña}; el PO \textbf{compartirá una cuenta de correo} para la configuración.} & \vl{Cambio} & \vl{Alta} & \vl{Cuenta de correo compartida por el PO (Oscar) y configuración de envío} & \vl{\textbf{Pendiente}} \\ \hline
REQ-27 & \vl{\textbf{Presentación/entrega del Análisis y Diseño (A\&D) por correo} al cliente.} & \vl{Documento} & \vl{Alta} & \vl{Envío del A\&D por correo} & \vl{\textbf{En proceso}} \\ \hline
\end{longtable}}

% ============================ 7. ACUERDOS =================================
\asec{7. ACUERDOS Y COMPROMISOS}
\vspace{5pt}
{\footnotesize
\begin{longtable}{|C{0.8cm}|L{5.6cm}|L{2.5cm}|C{1.8cm}|C{1.5cm}|L{2.0cm}|}
\hline
\rowcolor{grisH}\hc{N°} & \hc{Acción / entregable} & \hc{Responsable} & \hc{Fecha límite} & \hc{Estado} & \hc{Evidencia} \\ \hline
\endfirsthead
\hline
\rowcolor{grisH}\hc{N°} & \hc{Acción / entregable} & \hc{Responsable} & \hc{Fecha límite} & \hc{Estado} & \hc{Evidencia} \\ \hline
\endhead
1 & \vl{\textbf{Mostrar el avance desplegado y el módulo de autenticación} en los entornos de la universidad} & \vl{Juan Marcelo (Backend)} & \vl{21/09/2026 (en reunión)} & \vl{\textbf{Realizado}} & \vl{Esta acta / demo} \\ \hline
2 & \vl{\textbf{Compartir una cuenta de correo} para la configuración de envío con Gmail} & \vl{Oscar (Cliente)} & \vl{Por definir} & \vl{\textbf{Pendiente}} & \vl{Cuenta de correo} \\ \hline
3 & \vl{\textbf{Conseguir la API/configuración de Gmail (SMTP)} y configurar el envío de correos (recuperación de contraseña)} & \vl{Juan Marcelo (Backend)} & \vl{Por definir} & \vl{\textbf{Pendiente}} & \vl{Configuración Gmail} \\ \hline
4 & \vl{\textbf{Presentar/entregar el Análisis y Diseño (A\&D) por correo}} & \vl{Fabricio (PM) / Luis / Carlos (Analistas)} & \vl{22/09/2026} & \vl{\textbf{En proceso}} & \vl{Correo A\&D} \\ \hline
5 & \vl{\textbf{Enviar el acta 7 por correo}} & \vl{Fabricio (PM)} & \vl{22/09/2026} & \vl{\textbf{Realizado}} & \vl{Correo} \\ \hline
\end{longtable}}

% ============================ 8. PRÓXIMA ==================================
\asec{8. PRÓXIMA INTERACCIÓN Y SEGUIMIENTO}
{\footnotesize
\begin{tabular}{|L{2.5cm}|L{4.6cm}|L{2.6cm}|L{4.5cm}|}
\hline
\lb{Próxima reunión} & \vl{\textbf{Por definir} --- cadencia de \textbf{2 reuniones por semana}} & \lb{Objetivo} & \vl{Validar el despliegue y la configuración del envío de correos con Gmail, y registrar la \textbf{conformidad del PO} sobre el \textbf{Análisis y Diseño (A\&D)} remitido hoy por correo (cierre de la revisión)} \\ \hline
\lb{Canal de seguimiento} & \vl{WhatsApp (coordinación diaria) / Correo (actas)} & \lb{Responsable del acta} & \vl{Fabricio Godofredo Ladera La Torre (PM)} \\ \hline
\lb{Fecha de envío del acta} & \vl{22/09/2026} & \lb{Ubicación del archivo} & \vl{Repositorio del equipo CS3081} \\ \hline
\end{tabular}}

% ============================ 9. VALIDACIÓN ===============================
\asec{9. VALIDACIÓN DEL ACTA/FIRMAS}
{\footnotesize\itshape\color{grisN} La validación puede realizarse mediante firma, correo, mensaje en el canal acordado o confirmación registrada. Adjunte la evidencia correspondiente.\par}
\vspace{5pt}
{\footnotesize
\begin{tabular}{|L{4.7cm}|L{4.7cm}|L{4.7cm}|}
\hline
\hc{Representante del usuario / cliente} & \hc{Representante del equipo} & \hc{Docente o ACL (si corresponde)} \\ \hline
\vl{Nombre: Oscar Rafael Baldeón Montoro\newline Cargo / rol: Cliente / Product Owner} & \vl{Nombre: Fabricio Godofredo Ladera La Torre\newline Cargo / rol: Jefe de Proyecto (PM)} & \vl{Nombre: Teofilo Chambilla Aquino\newline Cargo / rol: Docente (acompañamiento)} \\ \hline
\vl{Método de validación: Firma y correo} & \vl{Método de validación: Correo (envío del acta)} & \vl{Método de validación: [Firma / correo]} \\ \hline
\vl{Fecha: 21/09/2026\newline Observaciones: Conforme con el avance desplegado y la configuración de correo} & \vl{Fecha: 21/09/2026\newline Observaciones: Acta enviada por correo} & \vl{Fecha: [dd/mm/aaaa]\newline Observaciones: [Completar]} \\ \hline
\end{tabular}}

% ============================ 10. APROBACIÓN ==============================
\asec{10. APROBACIÓN (V°B° CEO)}
El presente acta es revisada y aprobada por la Gerencia General antes de su envío al cliente.
\vspace{4pt}

{\itshape Oscar Rafael Baldeón Montoro}
\vspace{10pt}

\textbf{Nombre:} Oscar Rafael Baldeón Montoro

\textbf{Cargo:} CEO / Gerente General

\vspace{28pt}

\textbf{Firma:} \rule[-2pt]{6.5cm}{0.4pt}

\textbf{Fecha:} [dd/mm/aaaa]

\vspace{6pt}
\label{LastPage}
\end{document}
